\chapter*{Preface}
\addcontentsline{toc}{chapter}{Preface}

RS3D is an application for the Raspberry Shake 3D seismograph. Its \emph{User's Guide} explains how to use it;
this manual explains \emph{why it works}. It documents, one function at a time, the signal-processing theory that the
application relies on, how that theory was turned into code, and where the code deliberately or unavoidably departs
from the textbook version.

The manual is written for three kinds of reader: an observer who wants to trust (or question) what a spectrogram or an
alert is telling them, a student who wants a worked example of applied seismological signal processing, and a developer
who is about to change one of the algorithms and needs the assumptions written down.

\paragraph{What is new here, and what is not.}
Nothing in this manual is new science. Every method (Butterworth filtering, Welch spectra, short-time Fourier analysis,
the STA/LTA trigger) is standard and is cited. What the manual adds is precision about \emph{this} implementation:
the exact windows, constants, tolerances and edge-case rules, each traceable to a source file of the application.

\paragraph{How the chapters are organised.}
Chapters follow the path of a sample through the application, from the sensor (Chapters~\ref{ch:data}
and~\ref{ch:response}) through filtering and spectral analysis (Chapters~\ref{ch:filter}--\ref{ch:noise}) and data
quality (Chapter~\ref{ch:quality}) to visualisation, detection and real-time handling
(Chapters~\ref{ch:vis}--\ref{ch:realtime}). Chapter~\ref{ch:arch} describes how the code is organised, and the
appendices collect derivations and every numerical constant used by the application.

\begin{rs3dkey}[How to read the boxes]
\textbf{Key idea} boxes state a result in one or two sentences.\par\smallskip
\textbf{In RS3D} boxes say exactly what the application does, naming the source file (paths are relative to
\file{rs3d-app/src/rs3d/}). Where the application departs from the textbook method this is stated there, never
left implicit.
\end{rs3dkey}

\paragraph{Reproducibility.}
Every plot is produced by one script, \file{make\_figures.py}, which calls the application's own
functions (not re-implementations) on a seeded synthetic recording of one day at a fictitious station
(\file{AM.DEMO1.00}). The recording contains a local earthquake, a distant earthquake, two man-made noise bursts and two
data gaps, so each algorithm can be shown on a case whose true answer is known.

\paragraph{Status of the alert chapter.}
The event-detection theory of Chapter~\ref{ch:detect} is complete. The design of the application's alert pipeline in
Section~\ref{sec:alertdesign} is a \emph{design proposal}: that feature is not yet implemented, and the proposal
will be revised when it is.

\paragraph{Verification of references.}
Two of the key references (Trnkoczy 2012 and McNamara and Buland 2004) were checked against online bibliographic
records while this edition was written. The remaining entries are well-known works cited from the authors'
knowledge; their volume and page numbers should be confirmed against the originals before the references are
reused in a formal publication. The description of rsudp's alert defaults in Chapter~\ref{ch:detect} is deliberately
hedged for the same reason: no authoritative list of defaults could be located.

\bigskip
\begin{center}\small\textsf{\textcolor{rs3dgrey}{RS3D project, October 2026}}\end{center}

\chapter*{Notation}
\addcontentsline{toc}{chapter}{Notation}
\begin{center}
\small
\begin{tabularx}{\linewidth}{@{}l X l@{}}
\toprule
Symbol & Meaning & Unit\\ \midrule
$x(t),\;x[n]$ & continuous signal; its samples $x[n]=x(n\,\dt)$ & counts or m/s\\
$\fs,\;\dt$ & sampling frequency, sampling interval $\dt = 1/\fs$ & Hz, s\\
$\fN$ & Nyquist frequency, $\fN=\fs/2$ & Hz\\
$N$ & number of samples in a window or record & --\\
$T$ & duration $T=N\,\dt$ & s\\
$X[k]$ & discrete Fourier transform of $x[n]$ & counts\\
$f_k$ & frequency of bin $k$, $f_k = k\,\fs/N$ & Hz\\
$\df$ & frequency resolution (bin spacing) $\df=\fs/N=1/T$ & Hz\\
$w[n]$ & window (taper) function & --\\
$H(f)$ & transfer function of a filter & --\\
$T(f)$ & instrument response, counts per unit of ground velocity & counts/(m/s)\\
$S_{xx}(f)$ & one-sided power spectral density & counts$^2$/Hz\\
$\mathrm{STA},\mathrm{LTA}$ & short- and long-term averages & counts$^2$\\
$R$ & STA/LTA ratio & --\\
$\jj$ & imaginary unit & --\\
\bottomrule
\end{tabularx}
\end{center}
